% ---------------------------------------------------------------------------
% Drop-in replacement for Appendix app:constants of paper_new.tex
% (replaces lines 3469-3559 of the current file).
%
% All labels used elsewhere in the paper (app:constants, eq:app-caseA,
% eq:app-caseB) are preserved.  The scripted sweep over odd 63<=m<=499 and the
% twelve scripted integer comparisons are eliminated: every step below is
% analytic, and every rational constant is an exact fraction rounded in the
% safe direction, with the underlying comparison displayed.
%
% Every displayed inequality of this appendix is checked by
% validate_app_constants.py (dense float grids over the exact stated domains
% plus exact Fraction arithmetic at all corners and for all rational
% identities/comparisons).
% ---------------------------------------------------------------------------

\section{Constant inequalities for the residual-strip closure}\label{app:constants}

This appendix proves the two numerical inequalities used in
Lemma~\ref{lem:left-surplus-tail}.  Let
\[
        P(\theta)=\frac{2/3-2\theta}{\theta(1/2-2\theta)^2},
\]
\[
        B_0(\theta)=
        \left(\frac76\right)^{1-2\theta}
        (8-24\theta)^{-\theta}
        \frac{1/2+\theta}{5/12+\theta},
\]
and
\[
        B_1(\theta)=
        \left(\frac76\right)^{1-2\theta}
        (8-24\theta)^{-\theta}
        \frac{34}{29}.
\]
The required inequalities are
\begin{align}
        \frac{99}{100m}P(\theta)B_0(\theta)^m&\ge1
        &&(\theta=r/m\le1/6),\label{eq:app-caseA}\\
        \frac{99}{100m}P(\theta)B_1(\theta)^m&\ge1
        &&(1/6\le\theta=r/m<1/4),\label{eq:app-caseB}
\end{align}
for the integer pairs appearing in the residual strip with $m\ge63$ and
$r\ge2$.  We prove the following stronger statements, in which the parameters
are \emph{real} and unconstrained: \eqref{eq:app-caseA} holds for every real
$0<\theta\le1/6$ and every real $m>0$
(Proposition~\ref{prop:app-caseA-full}), and \eqref{eq:app-caseB} holds for
every real $1/6\le\theta<1/4$ and every real $m>0$
(Proposition~\ref{prop:app-caseB-full}).  In particular no finite
verification, and no use of the integrality of $r$ or of the parity or size
of $m$, remains in this appendix.

The elimination of the parameter $m$ rests on the elementary bound
$B^m/m\ge e\log B$ (Lemma~\ref{lem:app-tools}(ii) below): it replaces the
minimisation of $B^m/m$ over $m$ by its exact infimum over real $m>0$, which
is attained as $m\log B\to1$ and equals $e\log B$.  After this step both
inequalities become one-variable statements about $\theta$ alone, and these
are proved by exhibiting rational tangent/secant minorants of the logarithms
involved.

\subsection*{Elementary tools}

\begin{lemma}[Exponential and logarithm bounds]\label{lem:app-tools}
\begin{enumerate}[label=\textup{(\roman*)}]
\item $e^{y}\ge1+y$ for every real $y$.
\item For every real $B>0$ and every real $m>0$,
\[
        \frac{B^m}{m}\ \ge\ e\log B .
\]
\item $\log u\le u-1$ for every real $u>0$.
\item For every real $t\ge1$,
\[
        \frac{3(t^2-1)}{t^2+4t+1}
        \ \le\ \log t\ \le\
        \frac{(t-1)(t+5)}{2(2t+1)} .
\]
We write $L(t)$ and $U(t)$ for the left and right members.
\item $e\ge 65/24$.
\end{enumerate}
\end{lemma}

\begin{proof}
(i) The function $y\mapsto e^y-1-y$ has derivative $e^y-1$, which is
negative for $y<0$ and positive for $y>0$; its minimum value, at $y=0$, is
$0$.

(ii) Apply (i) with $y=m\log B-1$:
\[
        B^m=e^{m\log B}=e\cdot e^{\,m\log B-1}\ge e\,(1+m\log B-1)=e\,m\log B,
\]
and divide by $m>0$.

(iii) Apply (i) with $y=u-1$: $e^{u-1}\ge u$, and take logarithms (both
sides are positive).

(iv) Both bounds are equalities at $t=1$, so it suffices to compare
derivatives for $t\ge1$.  For the lower bound,
\[
        \frac{d}{dt}L(t)
        =\frac{6t(t^2+4t+1)-3(t^2-1)(2t+4)}{(t^2+4t+1)^2}
        =\frac{12(t^2+t+1)}{(t^2+4t+1)^2},
\]
and therefore
\[
        \frac{1}{t}-\frac{d}{dt}L(t)
        =\frac{(t^2+4t+1)^2-12t(t^2+t+1)}{t\,(t^2+4t+1)^2}
        =\frac{t^4-4t^3+6t^2-4t+1}{t\,(t^2+4t+1)^2}
        =\frac{(t-1)^4}{t\,(t^2+4t+1)^2}\ \ge\ 0 .
\]
For the upper bound, writing $U(t)=(t^2+4t-5)/(4t+2)$,
\[
        \frac{d}{dt}U(t)
        =\frac{(2t+4)(4t+2)-4(t^2+4t-5)}{(4t+2)^2}
        =\frac{4t^2+4t+28}{(4t+2)^2},
\]
and therefore
\[
        \frac{d}{dt}U(t)-\frac1t
        =\frac{4t^3+4t^2+28t-(4t+2)^2}{t\,(4t+2)^2}
        =\frac{4(t-1)^3}{t\,(4t+2)^2}\ \ge\ 0
        \qquad(t\ge1).
\]
Integrating the two derivative inequalities from $1$ to $t$ gives (iv).

(v) $e=\sum_{k\ge0}1/k!\ge1+1+\frac12+\frac16+\frac1{24}=\frac{65}{24}$,
since all terms of the series are positive.
\end{proof}

We record the six rational values of $L$ and $U$ that will be used; each is
obtained from Lemma~\ref{lem:app-tools}(iv) by clearing denominators.
\begin{equation}\label{eq:app-pade-values}
\begin{aligned}
        \log\frac{64}{63}&\ \ge\ L\!\left(\frac{64}{63}\right)
        =\frac{3\,(64^2-63^2)}{64^2+4\cdot64\cdot63+63^2}
        =\frac{3\cdot127}{24193},
        &
        \log\frac73&\ \ge\ L\!\left(\frac73\right)
        =\frac{3\,(49-9)}{49+84+9}
        =\frac{60}{71},
        \\[1mm]
        \log\frac{34}{29}&\ \ge\ L\!\left(\frac{34}{29}\right)
        =\frac{3\,(34^2-29^2)}{34^2+4\cdot34\cdot29+29^2}
        =\frac{945}{5941},
        &
        \log\frac76&\ \ge\ L\!\left(\frac76\right)
        =\frac{3\,(49-36)}{49+168+36}
        =\frac{39}{253},
        \\[1mm]
        \log 2&\ \le\ U(2)
        =\frac{1\cdot7}{2\cdot5}
        =\frac{7}{10},
        &
        \log\frac76&\ \le\ U\!\left(\frac76\right)
        =\frac{\frac16\cdot\frac{37}6}{2\cdot\frac{10}3}
        =\frac{37}{240}.
\end{aligned}
\end{equation}

\subsection*{Case A: $0<\theta\le1/6$}

Throughout this subsection set
\[
        \lambda_*=\frac{127}{24193},
        \qquad
        s=\frac{89}{100}.
\]

\begin{lemma}[Linear minorant of $\log B_0$]\label{lem:app-B0-linear}
For every $0\le\theta\le\frac16$,
\[
        \log B_0(\theta)\ \ge\ \lambda_*+s\left(\frac16-\theta\right).
\]
\end{lemma}

\begin{proof}
\emph{Value at the right endpoint.}
Since $4^{-1/6}=2^{-1/3}$,
\[
        B_0(1/6)
        =\left(\frac76\right)^{2/3}4^{-1/6}\,\frac{2/3}{7/12}
        =7^{2/3}6^{-2/3}\cdot2^{-1/3}\cdot\frac{8}{7}
        =8\,(7\cdot6^2\cdot2)^{-1/3}
        =\frac{8}{504^{1/3}}
        =\frac{4}{63^{1/3}},
\]
so that $B_0(1/6)^3=64/63$ and, by \eqref{eq:app-pade-values},
\begin{equation}\label{eq:app-lam-star}
        \log B_0(1/6)=\frac13\log\frac{64}{63}
        \ \ge\ \frac13\cdot\frac{3\cdot127}{24193}
        =\frac{127}{24193}=\lambda_* .
\end{equation}

\emph{Slope.}  From
\[
        \log B_0(\theta)
        =(1-2\theta)\log\frac76-\theta\log(8-24\theta)
        +\log\Bigl(\frac12+\theta\Bigr)-\log\Bigl(\frac5{12}+\theta\Bigr),
\]
the negated derivative $c(\theta):=-\dfrac{d}{d\theta}\log B_0(\theta)$ is
\begin{equation}\label{eq:app-c-def}
        c(\theta)
        =2\log\frac76+\log(8-24\theta)-\frac{24\theta}{8-24\theta}
        -\frac{1}{1/2+\theta}+\frac{1}{5/12+\theta},
\end{equation}
and differentiating once more,
\begin{equation}\label{eq:app-c-prime}
        c'(\theta)
        =-\frac{24}{8-24\theta}-\frac{192}{(8-24\theta)^2}
        +\frac{1}{(1/2+\theta)^2}-\frac{1}{(5/12+\theta)^2}.
\end{equation}
On $[0,1/6]$ we have $8-24\theta\ge4>0$, so the first two terms of
\eqref{eq:app-c-prime} are negative; and $0<\frac5{12}+\theta<\frac12+\theta$
gives $(1/2+\theta)^{-2}<(5/12+\theta)^{-2}$, so the sum of the last two
terms is negative as well.  Hence $c'<0$ on $[0,1/6]$, i.e.\ $c$ is strictly
decreasing there, and for every $\theta\in[0,1/6]$,
\begin{equation}\label{eq:app-c-lb}
        c(\theta)\ \ge\ c(1/6)
        =2\log\frac76+\log4-1-\frac32+\frac{12}7
        =2\log\frac73-\frac{11}{14}
        \ \ge\ \frac{120}{71}-\frac{11}{14}
        =\frac{899}{994}
        \ \ge\ \frac{89}{100}=s,
\end{equation}
where the value of $c(1/6)$ is read off from \eqref{eq:app-c-def} (using
$2\log\frac76+\log4=2\log\frac73$ and $-1-\frac32+\frac{12}7=-\frac{11}{14}$),
the penultimate inequality is \eqref{eq:app-pade-values} at $t=7/3$, and the
last comparison is $899\cdot100=89900\ge88466=89\cdot994$.

\emph{Conclusion.}  By the fundamental theorem of calculus, for
$0\le\theta\le1/6$,
\[
        \log B_0(\theta)
        =\log B_0(1/6)+\int_\theta^{1/6}c(u)\,du
        \ \ge\ \lambda_*+s\left(\frac16-\theta\right). \qedhere
\]
\end{proof}

\begin{lemma}[Weighted polynomial inequality]\label{lem:app-cubic}
For every $0<\theta\le\frac16$,
\[
        P(\theta)\left[\lambda_*+s\left(\frac16-\theta\right)\right]
        \ \ge\ 72\,\lambda_* .
\]
\end{lemma}

\begin{proof}
For $0<\theta<1/4$ one has $\theta(1/2-2\theta)^2>0$, so the claim is
equivalent to $G(\theta)\ge0$ on $(0,1/6]$, where
\[
        G(\theta)
        :=\left(\frac23-2\theta\right)
        \left[\lambda_*+s\left(\frac16-\theta\right)\right]
        -72\lambda_*\,\theta\left(\frac12-2\theta\right)^2 .
\]
$G$ is a cubic polynomial with rational coefficients.  At $\theta=1/6$ the
first term equals $\frac13\lambda_*$ and the second equals
$72\lambda_*\cdot\frac16\cdot\frac1{36}=\frac13\lambda_*$, so $G(1/6)=0$ and
$\left(\frac16-\theta\right)$ divides $G$.  Polynomial division gives the
exact factorisation
\begin{equation}\label{eq:app-G-factor}
        G(\theta)=\left(\frac16-\theta\right)Q(\theta),
        \qquad
        Q(\theta)
        =288\lambda_*\theta^2-(2s+96\lambda_*)\theta
        +\left(\frac{2s}3+4\lambda_*\right),
\end{equation}
which is verified by expanding both sides to
\[
        \left(\frac{2\lambda_*}3+\frac s9\right)
        -\left(s+20\lambda_*\right)\theta
        +\left(2s+144\lambda_*\right)\theta^2
        -288\lambda_*\theta^3 .
\]
On $[0,1/6]$,
\[
        Q'(\theta)=576\lambda_*\theta-(2s+96\lambda_*)
        \ \le\ 96\lambda_*-(2s+96\lambda_*)=-2s\ <\ 0,
\]
so $Q$ is strictly decreasing on $[0,1/6]$, and
\begin{equation}\label{eq:app-Q-endpoint}
        Q(\theta)\ \ge\ Q(1/6)
        =8\lambda_*-\frac{2s+96\lambda_*}{6}+\frac{2s}3+4\lambda_*
        =\frac s3-4\lambda_*
        =\frac{89}{300}-\frac{508}{24193}
        =\frac{2000777}{7257900}\ >\ 0 .
\end{equation}
Hence $Q>0$ on $[0,1/6]$ and $G\ge0$ on $(0,1/6]$ (with equality only at
$\theta=1/6$), which proves the lemma.
\end{proof}

\begin{proposition}[Case A, all real parameters]\label{prop:app-caseA-full}
For every real $0<\theta\le\frac16$ and every real $m>0$,
\[
        \frac{99}{100m}P(\theta)B_0(\theta)^m\ \ge\
        \frac{99}{100}\cdot72\cdot\frac{127}{24193}\cdot\frac{65}{24}
        =\frac{58841640}{58063200}\ >\ 1 .
\]
In particular \eqref{eq:app-caseA} holds for every integer pair of the
residual strip.
\end{proposition}

\begin{proof}
By Lemma~\ref{lem:app-B0-linear},
$\log B_0(\theta)\ge\lambda_*+s(\frac16-\theta)>0$.  By
Lemma~\ref{lem:app-tools}(ii) applied to $B=B_0(\theta)$,
\[
        \frac{B_0(\theta)^m}{m}\ \ge\ e\log B_0(\theta)
        \ \ge\ e\left[\lambda_*+s\left(\frac16-\theta\right)\right].
\]
Since $0<\theta\le1/6$ gives $\frac23-2\theta\ge\frac13>0$, we have
$P(\theta)>0$, and multiplying by $\frac{99}{100}P(\theta)$ preserves the
inequality:
\[
        \frac{99}{100m}P(\theta)B_0(\theta)^m
        \ \ge\ \frac{99e}{100}\,P(\theta)
        \left[\lambda_*+s\left(\frac16-\theta\right)\right]
        \ \ge\ \frac{99e}{100}\cdot72\,\lambda_*,
\]
the last step by Lemma~\ref{lem:app-cubic}.  Finally, by
Lemma~\ref{lem:app-tools}(v),
\[
        \frac{99e}{100}\cdot72\,\lambda_*
        \ \ge\ \frac{99}{100}\cdot72\cdot\frac{127}{24193}\cdot\frac{65}{24}
        =\frac{99\cdot72\cdot127\cdot65}{100\cdot24193\cdot24}
        =\frac{58841640}{58063200}\ >\ 1 . \qedhere
\]
\end{proof}

\subsection*{Case B: $1/6\le\theta<1/4$}

\begin{lemma}[Prefactor bound]\label{lem:app-P72}
For every $\frac16\le\theta<\frac14$, $\;P(\theta)\ge72$.
\end{lemma}

\begin{proof}
One checks by expansion the polynomial identity
\[
        \frac23-2\theta-72\,\theta\left(\frac12-2\theta\right)^2
        =-\frac23\,(6\theta-1)(72\theta^2-24\theta+1).
\]
On $[1/6,1/4]$ the factor $6\theta-1$ is nonnegative.  The quadratic
$72\theta^2-24\theta+1$ is convex, and its values at the endpoints are
\[
        72\cdot\frac1{36}-24\cdot\frac16+1=-1,
        \qquad
        72\cdot\frac1{16}-24\cdot\frac14+1=-\frac12,
\]
so by convexity it is at most $\max\{-1,-\frac12\}<0$ on all of $[1/6,1/4]$.
Hence the identity's right-hand side is nonnegative on $[1/6,1/4]$, i.e.
\[
        \frac23-2\theta\ \ge\ 72\,\theta\left(\frac12-2\theta\right)^2
        \qquad\left(\frac16\le\theta\le\frac14\right),
\]
and dividing by $\theta(1/2-2\theta)^2>0$ (legitimate for $\theta<1/4$)
gives $P(\theta)\ge72$.
\end{proof}

\begin{lemma}[Uniform lower bound for $\log B_1$]\label{lem:app-B1-quad}
For every $\frac16\le\theta\le\frac14$,
\[
        \log B_1(\theta)\ \ge\
        \frac{945}{5941}+\frac{39}{253}-\frac{4225}{13824}
        =\frac{157634591}{20778481152}
        \ \ge\ \frac{7}{1000}.
\]
\end{lemma}

\begin{proof}
Write $8-24\theta=4(2-6\theta)$, where $2-6\theta\in[\frac12,1]$ on the
stated range.  By Lemma~\ref{lem:app-tools}(iii) with $u=2-6\theta$,
\[
        \log(2-6\theta)\ \le\ 1-6\theta,
\]
hence, using $\theta>0$,
\[
        \log B_1(\theta)
        =\log\frac{34}{29}+(1-2\theta)\log\frac76
        -\theta\log4-\theta\log(2-6\theta)
        \ \ge\ C-K\theta+6\theta^2,
\]
where
\[
        C=\log\frac{34}{29}+\log\frac76,
        \qquad
        K=1+2\log2+2\log\frac76 .
\]
By \eqref{eq:app-pade-values},
\[
        C\ \ge\ \bar C:=\frac{945}{5941}+\frac{39}{253},
        \qquad
        K\ \le\ \bar K:=1+2\cdot\frac{7}{10}+2\cdot\frac{37}{240}
        =1+\frac75+\frac{37}{120}=\frac{325}{120}=\frac{65}{24}.
\]
Since $\theta>0$, replacing $C$ by $\bar C$ and $K$ by $\bar K$ can only
decrease $C-K\theta+6\theta^2$, and completing the square,
\[
        \log B_1(\theta)
        \ \ge\ \bar C-\bar K\theta+6\theta^2
        =\bar C-\frac{\bar K^2}{24}+6\left(\theta-\frac{\bar K}{12}\right)^2
        \ \ge\ \bar C-\frac{\bar K^2}{24}
        =\frac{945}{5941}+\frac{39}{253}-\frac{4225}{13824}.
\]
The exact value of the right-hand side is
$\frac{157634591}{20778481152}$ (the three denominators
$5941=13\cdot457$, $253=11\cdot23$ and $13824=2^9\cdot3^3$ are pairwise
coprime), and the final comparison
\[
        \frac{157634591}{20778481152}\ \ge\ \frac{7}{1000}
\]
is the integer inequality
$1000\cdot157634591=157634591000\ \ge\ 145449368064=7\cdot20778481152$.
\end{proof}

\begin{proposition}[Case B, all real parameters]\label{prop:app-caseB-full}
For every real $\frac16\le\theta<\frac14$ and every real $m>0$,
\[
        \frac{99}{100m}P(\theta)B_1(\theta)^m\ \ge\
        \frac{99}{100}\cdot72\cdot\frac{7}{1000}\cdot\frac{65}{24}
        =\frac{3243240}{2400000}\ >\ 1 .
\]
In particular \eqref{eq:app-caseB} holds for every integer pair of the
residual strip.
\end{proposition}

\begin{proof}
By Lemma~\ref{lem:app-B1-quad}, $\log B_1(\theta)\ge\frac7{1000}>0$.  By
Lemma~\ref{lem:app-tools}(ii) applied to $B=B_1(\theta)$, and then
Lemma~\ref{lem:app-P72} together with $\log B_1(\theta)>0$,
\[
        \frac{99}{100m}P(\theta)B_1(\theta)^m
        \ \ge\ \frac{99e}{100}\,P(\theta)\log B_1(\theta)
        \ \ge\ \frac{99e}{100}\cdot72\cdot\frac{7}{1000}
        \ \ge\ \frac{99}{100}\cdot72\cdot\frac{7}{1000}\cdot\frac{65}{24},
\]
the last step by Lemma~\ref{lem:app-tools}(v).  Finally
\[
        \frac{99\cdot72\cdot7\cdot65}{100\cdot1000\cdot24}
        =\frac{3243240}{2400000}\ >\ 1 . \qedhere
\]
\end{proof}

\begin{remark}[No computer assistance]\label{rmk:app-no-computer}
Propositions~\ref{prop:app-caseA-full} and~\ref{prop:app-caseB-full} prove
\eqref{eq:app-caseA} and \eqref{eq:app-caseB} for all real $\theta$ in the
stated ranges and all real $m>0$; the hypotheses $m\ge63$, $r\ge2$, $n>2r$
and $\theta=r/m$ of the residual strip are not needed here.  Consequently
the scripted verification of the finite range odd $63\le m\le499$ and the
twelve scripted integer comparisons of earlier versions of this appendix are
now superfluous, and the only computer-assisted ingredient remaining in the
proof of Theorem~\ref{thm:regionI-full} is Proposition~\ref{prop:M61}
(the residual strip for $m\le61$).  The margins of the two propositions are
$\frac{58841640}{58063200}-1\approx1.3\%$ (Case A; this is intrinsically
tight, since $\frac{99}{100m}P(\theta)B_0(\theta)^m\to\frac{99e}{100}\cdot
72\log B_0(1/6)\approx1.017$ along $\theta=1/6$, $m\log B_0(1/6)\to1$) and
$\frac{3243240}{2400000}-1\approx35\%$ (Case B).
\end{remark}
